% Isolated exact-reduction candidate. No geometric nonvanishing is claimed.
\subsection{The precise restricted cup-product pairing}

\begin{lemma}[The actual determinant directions and their dual pairing]
\label{lem:DE-A3-positive-restricted-Petri-note}
\status{Candidate equivalent reduction; complete proof written, independent audit pending.}
Keep a sharp source in the preceding monad construction, put
$\mathcal L=\det Q=LP$, and write $U=H^1(T,Q^\vee)$.
The determinant-direction map has the exact sequence
\[
0\longrightarrow H^0(KP)\longrightarrow
H^0(E_0)\oplus H^0(E_1)\xrightarrow{\dot w}
H^0(\mathcal L)\longrightarrow H^1(KP)\longrightarrow0.
\tag{DE-PT1}
\]
In particular, its image $\Delta_s$ satisfies
\[
\operatorname{codim}\Delta_s=h^1(KP)=h^0(KP).
\tag{DE-PT2}
\]
The symmetric bilinear cup-wedge map
\[
\mu:\operatorname{Sym}^2 U\longrightarrow
H^2(T,\mathcal L^{-1}),\qquad
x\,y\longmapsto \wedge_*(x\mathbin{\smile}y)
\tag{DE-PT3}
\]
is well defined. Choose the fixed invariant volume form to identify
$H^2(T,\mathcal L^{-1})=H^0(T,\mathcal L)^*$ and let
$\mu_{\Delta}$ be its composite with restriction to $\Delta_s^*$.
Then $\mu_{\Delta}$ is the linear dual of the family
\[
\Psi_s:\Delta_s\longrightarrow
\operatorname{Hom}_k(U,H^1(Q)),\qquad
v\longmapsto H^1(vJ_Q),
\]
after Serre duality identifies $H^1(Q)$ with $U^*$.
Consequently
\[
\dim\operatorname{im}\Psi_s=\dim\operatorname{im}\mu_{\Delta}.
\tag{DE-PT4}
\]
No smoothness or reducedness of the determinant curve is required.
\end{lemma}
\begin{proof}
Twist the defining sequence for $K$ by $P$ and interchange its two
summands using $E_1=E_0P$ and the fixed $P^2$ trivialization.
Its map to $LP$ becomes
$(t_0,t_1)\mapsto s_0\wedge t_1+t_0\wedge s_1$.
The absence of a common zero of the two vector-bundle sections makes
this sheaf map surjective. Since $H^j(E_i)=0$ for $j>0$,
its cohomology gives (DE-PT1). Both the middle section space and
$H^0(LP)$ have dimension $b$, proving (DE-PT2).

Interchanging $x,y$ in degree-one cup product contributes its graded
sign, and interchanging the two factors of the alternating bundle
pairing contributes the same sign. Their product is positive.
Thus the bilinear map in (DE-PT3) is symmetric, including in
characteristic two, and factors through the stated symmetric-power
quotient. No identification of symmetric and divided powers is used.

For $v\in\Delta_s$ and $x,y\in U$, functoriality of cup product and
Serre duality gives
\[
\bigl\langle H^1(vJ_Q)x,y\bigr\rangle
=\operatorname{Tr}_T\bigl(v\,\wedge_*(x\mathbin{\smile}y)\bigr).
\tag{DE-PT5}
\]
This can also be checked in a frame of $Q$, where $J_Q$ is the
inverse alternating matrix and contraction with $y$ is the wedge
pairing on $Q^\vee$. Equation (DE-PT5) is precisely the duality
assertion. The image of $\Psi_s$ consists of symmetric bilinear
forms, which is the space $(\operatorname{Sym}^2U)^*$; taking the
dual map therefore proves (DE-PT4). All operations are on vector
bundles and their cohomology on the smooth surface $T$.
\end{proof}

\begin{corollary}[An exact common-radical test]
\label{cor:DE-A3-positive-Petri-radical-note}
\status{Candidate equivalent reduction; no radical vanishing asserted.}
At the fixed sharp source, a trace-direction variation with nonzero
determinant exists if and only if
\[
\{x\in U:\mu_{\Delta}(x\,y)=0\text{ for every }y\in U\}=0.
\tag{DE-PT6}
\]
Equivalently, no nonzero $x\in U$ satisfies
\[
\mu(xU)\subset\Delta_s^\perp
=\operatorname{im}\bigl(H^1(KP)^*\longrightarrow
H^0(LP)^*\bigr).
\tag{DE-PT7}
\]
\end{corollary}
\begin{proof}
By (DE-PT5) the displayed radical is exactly the common kernel of
all the maps $\Psi_s(v)$. An invertible member forces that common
kernel to be zero. Conversely, the proved trace-zero theorem says
that, in the compatible graph bases, a linear space consisting only
of singular variations is zero or a single nilpotent matrix line.
Such a space has a nonzero common kernel. Changing to the
Serre-duality bases preserves the common-kernel condition, so a zero
radical forces an invertible member. The final equality follows
by dualizing (DE-PT1).
\end{proof}

\begin{remark}[What is still missing]
\label{rem:DE-A3-positive-Petri-radical-scope-note}
\status{Unproved geometric gate.}
Neither (DE-PT2) nor (DE-PT3) bounds the restricted cup-product
image from below. In particular the determinant directions cannot
be replaced by all of $H^0(LP)$: their actual annihilator is the
dual of $H^1(KP)$, with the twist by $P$ essential.
The curve pencil trick, when its additional hypotheses hold, controls
a different multiplication map on the determinant curve; no
identification with (DE-PT3) is assumed. The remaining uniform
problem is to prove (DE-PT6), or a sufficient two-dimensional image
bound in (DE-PT4), at one suitable sharp source on each surface.
\end{remark}
